\documentclass[11pt,oneside,reqno]{amsbook}
\usepackage[utf8]{inputenc}

\usepackage{enumitem}

\usepackage{amsthm}
\usepackage{mathtools}
\usepackage{amsmath}
\usepackage{amssymb}
\usepackage[all]{xy}
\usepackage{tikz}
\usepackage{tikz-cd}
\usepackage{mathtools}
\usepackage{mathrsfs}

\usepackage{color}

\def\rcomment#1{{\color{red}[{\sf #1}]}}
\def\bcomment#1{{\color{blue}[{\sf #1}]}}
\def\gcomment#1{{\color{green}[{\sf #1}]}}

%% blackboard bold

\newcommand\Z{{\mathbb Z}}
\newcommand\Q{{\mathbb Q}}
\newcommand\Qbar{{\overline{\mathbb Q}}}
\newcommand\N{{\mathbb N}}
\newcommand\R{{\mathbb R}}
\newcommand\C{{\mathbb C}}

\DeclareMathOperator{\sgn}{\text{sgn}}

%% bold face macros

\newcommand\br{{\mathbf r}}
\newcommand\bt{{\mathbf t}}
\newcommand\bu{\mathbf{u}}
\newcommand\bv{\mathbf{v}}
\newcommand\bw{\mathbf{w}}
\newcommand\bx{\mathbf{x}}
\newcommand\by{\mathbf{y}}

%%%%%%%%%%%%%%%%% environments %%%%%%%%%%%%%%%%%%

\numberwithin{equation}{section}

\newtheorem{theorem}{Theorem}[section]
\newtheorem{lemma}[theorem]{Lemma}
\newtheorem{prop}[theorem]{Proposition}
\newtheorem{cor}[theorem]{Corollary}

\theoremstyle{definition}
\newtheorem{definition}[theorem]{Definition}
\newtheorem{example}[theorem]{Example}
\newtheorem{problem}{Problem}
\newtheorem{step}{Step}

\theoremstyle{remark}
\newtheorem{remark}[theorem]{Remark}
\newtheorem{notation}[theorem]{Notation}
\newtheorem{question}[theorem]{Question}
\newtheorem{conjecture}[theorem]{Conjecture}

\title{Limits and Continuity}

\begin{document}

\section*{Limits and Continuity}

\subsection*{1} Evaluate the following limits, or state that they do not exist.
\[
\lim_{x \to 3} \frac{x^2 - 9}{x - 3}, \quad
\lim_{x \to 0} \frac{\sin x}{x}, \quad
\lim_{x \to \infty} \frac{3x^2 - 2x + 1}{x^2 + 5}, \quad
\lim_{x \to 0^+} \frac{1}{x}, \quad
\lim_{x \to 2} \frac{x^2 - 4x + 4}{x^2 - 4}.
\]

\subsection*{2} Suppose $\displaystyle\lim_{x \to a} f(x) = L$ and $\displaystyle\lim_{x \to a} g(x) = M$. Using only the intuitive idea that limits respect algebraic operations, justify each of the following:
\begin{gather*}
\lim_{x \to a}[f(x) + g(x)] = L + M, \qquad \lim_{x \to a}[f(x)g(x)] = LM, \\
\lim_{x \to a}\frac{f(x)}{g(x)} = \frac{L}{M} \quad (M \neq 0), \qquad \lim_{x \to a} c \cdot f(x) = cL.
\end{gather*}
Use these rules to evaluate $\displaystyle\lim_{x \to 2}(3x^2 - 5x + 1)$ and $\displaystyle\lim_{x \to 1} \frac{x^3 + 2x}{x^2 + 1}$.

\subsection*{3} Suppose $f(x) \leq g(x) \leq h(x)$ for all $x$ near $a$, and suppose
\[
\lim_{x \to a} f(x) = \lim_{x \to a} h(x) = L.
\]
\begin{enumerate}[label=(\alph*)]
    \item Explain informally why $\displaystyle\lim_{x \to a} g(x) = L$ must hold. This is called the \textit{Squeeze Theorem}.
    \item Use the Squeeze Theorem to evaluate $\displaystyle\lim_{x \to 0} x^2 \sin\!\left(\frac{1}{x}\right)$.
    \textit{Hint: what can you say about $|\sin(1/x)|$ for all $x \neq 0$?}
    \item Use the Squeeze Theorem and the geometric inequality
    \[
    \cos\theta \leq \frac{\sin\theta}{\theta} \leq 1, \qquad \theta \in \left(0, \tfrac{\pi}{2}\right),
    \]
    to show that $\displaystyle\lim_{\theta \to 0} \frac{\sin\theta}{\theta} = 1$.
\end{enumerate}

\subsection*{4} \textbf{Definition.} A function $f$ is \textit{continuous at} $a$ if
\[
\lim_{x \to a} f(x) = f(a).
\]
Note this requires three things: $f(a)$ is defined, $\lim_{x\to a}f(x)$ exists, and the two are equal.

For each function below, determine all points of discontinuity and classify each as \textit{removable}, \textit{jump}, or \textit{infinite}.
\[
f(x) = \frac{x^2-1}{x-1}, \quad g(x) = \begin{cases} x+1 & x < 0 \\ x^2 & x \geq 0 \end{cases}, \quad h(x) = \frac{1}{(x-2)^2}, \quad k(x) = \begin{cases} \dfrac{\sin x}{x} & x \neq 0 \\ 0 & x = 0 \end{cases}.
\]

\subsection*{5} Suppose $f$ and $g$ are continuous at $a$.
\begin{enumerate}[label=(\alph*)]
    \item Show that $f + g$ and $f \cdot g$ are continuous at $a$.
    \item Is $f/g$ necessarily continuous at $a$? If not, give a counterexample and state a condition under which it is.
    \item Suppose additionally that $g$ is continuous at $f(a)$. Show that the composition $g \circ f$ is continuous at $a$.
\end{enumerate}

\subsection*{6} Let $f(x) = \dfrac{|x|}{x}$ for $x \neq 0$.
\begin{enumerate}[label=(\alph*)]
    \item Compute $\displaystyle\lim_{x \to 0^+} f(x)$ and $\displaystyle\lim_{x \to 0^-} f(x)$.
    \item Explain why $\displaystyle\lim_{x \to 0} f(x)$ does not exist.
    \item Show that $f$ is continuous at every $a \neq 0$. \textit{Hint: consider $a > 0$ and $a < 0$ separately.}
    \item Can $f$ be defined at $x = 0$ in a way that makes it continuous there? Explain.
\end{enumerate}

\subsection*{7} \textbf{(Intermediate Value Theorem.)} Let $f$ be continuous on a closed interval $[a, b]$, and let $N$ be any value strictly between $f(a)$ and $f(b)$. Then there exists some $c \in (a,b)$ such that $f(c) = N$.
\begin{enumerate}[label=(\alph*)]
    \item Use the Intermediate Value Theorem to show that $x^5 - 3x + 1 = 0$ has at least one root in $[1, 2]$ and at least one root in $[-2, -1]$.
    \item Show that the equation $\cos x = x$ has at least one real solution.
    \textit{Hint: consider $f(x) = \cos x - x$ on a suitable interval.}
    \item Explain why every polynomial of odd degree with real coefficients must have at least one real root.
    \textit{Hint: think about the behavior of odd-degree polynomials as $x \to \pm\infty$.}
\end{enumerate}

\subsection*{8} Evaluate the following limits.
\[
\lim_{x \to \infty} \frac{5x^3 - x}{2x^3 + 7x^2 - 1}, \qquad
\lim_{x \to \infty}\!\left(\sqrt{x^2+x} - x\right), \qquad
\lim_{x \to -\infty} \frac{e^x}{1 + e^x}.
\]
For the middle limit, multiply and divide by the conjugate $\sqrt{x^2+x}+x$ and simplify before taking the limit.

\subsection*{9} Recall the difference quotient $A(f, h, x) = \dfrac{f(x+h)-f(x)}{h}$ from the Functions sheet.
\begin{enumerate}[label=(\alph*)]
    \item Compute $A(f, h, 4)$ for $f(x) = \sqrt{x}$ and simplify fully. Then evaluate $\displaystyle\lim_{h \to 0} A(f, h, 4)$.
    \item Repeat for $f(x) = \dfrac{1}{x}$ at general $x \neq 0$.
    \item For which value(s) of $x$ does the limit in part (b) fail to exist? What does this say about the graph of $f(x) = 1/x$?
    \item Compute $\displaystyle\lim_{h \to 0} A(f, h, x)$ for $f(x) = x^n$, $n$ a positive integer.
    \textit{Hint: use the factorization $a^n - b^n = (a-b)(a^{n-1}+a^{n-2}b + \cdots + b^{n-1})$ with $a = x+h$, $b = x$.}
\end{enumerate}

\subsection*{10} Suppose $f$ is continuous on $[a,b]$.
\begin{enumerate}[label=(\alph*)]
    \item Show that if $f(x) \neq 0$ for all $x \in [a,b]$, then $f$ must be either entirely positive or entirely negative on $[a,b]$. \textit{Hint: use the Intermediate Value Theorem.}
    \item Suppose $f(a) < 0 < f(b)$. Show that the set of roots $\{x \in [a,b] : f(x) = 0\}$ is nonempty. Does it have to be finite? Give an example where it is infinite.
    \item Give an example of a function on $[0,1]$ that is \textit{not} continuous but still satisfies: for every $y_0$ strictly between $f(0)$ and $f(1)$, there exists $c \in (0,1)$ with $f(c) = y_0$. What does this show about the converse of the IVT?
\end{enumerate}

\subsection*{11} Let $f(x) = \sin x$.
\begin{enumerate}[label=(\alph*)]
    \item Using the identity
    \[
    \sin x - \sin y = 2\cos\!\left(\frac{x+y}{2}\right)\sin\!\left(\frac{x-y}{2}\right),
    \]
    show that $|\sin x - \sin y| \leq |x - y|$ for all $x, y \in \R$.
    \item Conclude from part (a) that $\sin x$ is continuous at every $a \in \R$.
    \item Using the analogous identity for cosine, show that $\cos x$ is also continuous everywhere.
    \item Using parts (b) and (c) together with Problem~5, determine at which values of $x$ the function $\tan x = \sin x / \cos x$ is continuous.
\end{enumerate}

\end{document}
